\documentclass[12pt,a4paper]{article}
\usepackage[a4paper,margin=2cm]{geometry}
\usepackage[T2A]{fontenc}
\usepackage[utf8]{inputenc}
\usepackage[russian]{babel}
\usepackage{amsmath,amssymb,mathtools}
\usepackage{array,booktabs,longtable}
\usepackage{xcolor}
\usepackage{hyperref}
\usepackage{tikz}
\hypersetup{colorlinks=true,linkcolor=blue!50!black,urlcolor=blue!50!black}
\setlength{\parindent}{0pt}
\setlength{\parskip}{0.55em}
\setlength{\tabcolsep}{3pt}
\renewcommand{\arraystretch}{1.18}
\newcommand{\taskhead}[1]{\subsection*{#1}}
\newcommand{\field}[1]{\textbf{#1}}

\begin{document}

\begin{titlepage}
\thispagestyle{empty}
\begin{center}
\vspace*{0.22\textheight}
{\Huge\bfseries Ревью домашней работы\par}
\vspace{2.2cm}
{\Large Автор: Лёвин Артём Александрович\par}
\vspace{0.8cm}
{\Large 16 сентября 2026 года\par}
\vfill
\begin{tikzpicture}
\draw[blue!45!black,line width=0.8pt] (0,0) .. controls (2,0.65) and (4,-0.65) .. (6,0);
\end{tikzpicture}
\end{center}
\end{titlepage}

\newpage
\section*{Сводная таблица}
\small
\begin{longtable}{>{\raggedright\arraybackslash}p{0.13\linewidth} >{\raggedright\arraybackslash}p{0.19\linewidth} >{\raggedright\arraybackslash}p{0.26\linewidth} >{\raggedright\arraybackslash}p{0.15\linewidth} >{\raggedright\arraybackslash}p{0.17\linewidth}}
\toprule
\textbf{№ задания} & \textbf{Тема} & \textbf{Суть ошибки} & \textbf{Ваш ответ} & \textbf{Правильный ответ}\\
\midrule
\endfirsthead
\toprule
\textbf{№ задания} & \textbf{Тема} & \textbf{Суть ошибки} & \textbf{Ваш ответ} & \textbf{Правильный ответ}\\
\midrule
\endhead
\hyperref[task:14-2]{14.09 №2} & Равноускоренное движение & Скорость найдена верно, но не указано обязательное направление движения при положительной скорости. & $v_x=2\,\text{м/с}$ & $v_x=2\,\text{м/с}$, движение в положительном направлении оси\\
\midrule
\hyperref[task:14-3]{14.09 №3} & Формула скорости & Потерян отрицательный знак ускорения, поэтому начальная скорость вычислена неверно. & $v_0=2\,\text{м/с}$ & $v_0=26\,\text{м/с}$\\
\midrule
\hyperref[task:14-4]{14.09 №4} & Закон скорости и направление движения & Числовые характеристики найдены, но не описано движение до остановки и после неё. & $v_{0x}=7\,\text{м/с}$, $a_x=-1\,\text{м/с}^2$, $t_{\text{ост}}=7\,\text{с}$ & Те же значения; до $7$ с движение по положительному направлению, после $7$ с --- по отрицательному\\
\midrule
\hyperref[task:14-5]{14.09 №5} & Линейный закон скорости по двум точкам & Решение не доведено до вычисления ускорения, начальной скорости и закона $v_x(t)$; для двух произвольных точек нужно делить изменение скорости на изменение времени. & Не получен & $a_x=-2\,\text{м/с}^2$, $v_{0x}=8\,\text{м/с}$, $v_x=8-2t$\\
\midrule
\hyperref[task:15-1]{15.09 №1} & Ускорение по изменению скорости & В исходных данных потерян знак минус у начальной скорости, поэтому изменение скорости уменьшено с двенадцати до четырёх метров в секунду. & $a=\frac{2}{3}\,\text{м/с}^2$ & $a=2\,\text{м/с}^2$\\
\bottomrule
\end{longtable}
\normalsize

\newpage
\thispagestyle{empty}
\vspace*{\fill}
\begin{center}
{\Huge\bfseries Разбор ошибок}
\end{center}
\vspace*{\fill}
\newpage

\phantomsection\label{task:14-2}
\taskhead{14.09.2026 --- задание №2}

\field{Текст задания.} Тело имело начальную скорость $-6\,\text{м/с}$ и ускорение $2\,\text{м/с}^2$. Найдите скорость через $4$ с и укажите направление движения в этот момент.

\field{Ваше решение.}
\[
v_x=-6+2\cdot4=2\,\text{м/с}.
\]

\field{Ваш ответ.} $v_x=2\,\text{м/с}$. Направление движения отдельно не указано.

\textcolor{red}{\textbf{Ошибка:}} пропущен обязательный элемент ответа --- направление движения.

\field{Где именно возникает ошибка.} Последний правильный шаг --- вычисление скорости: $v_x=2\,\text{м/с}$. Первый неполный шаг --- запись окончательного ответа без интерпретации знака найденной скорости.

\field{Почему это неверно.} В условии требуется не только числовое значение скорости, но и направление движения. В одномерном движении знак проекции скорости определяет направление относительно выбранной положительной оси. Положительное значение $v_x$ означает движение в положительном направлении оси; отрицательное --- в противоположном. Поэтому ответ без указания направления не полностью выполняет условие, даже если вычисление скорости верно.

\textcolor{blue!60!black}{\textbf{Ключевая идея:}} после вычисления проекции скорости всегда нужно прочитать её знак и перевести его в словесное описание направления движения, если это прямо требуется условием.

\field{Исправление.} Полученное значение $v_x=2\,\text{м/с}$ сохраняем. Поскольку оно положительное, в момент $t=4$ с тело движется в положительном направлении оси.

\field{Полное правильное решение.} Для равноускоренного движения
\[
v_x=v_{0x}+a_xt.
\]
Подставляем данные:
\[
v_x=-6+2\cdot4=-6+8=2\,\text{м/с}.
\]
Так как $v_x>0$, направление движения совпадает с положительным направлением выбранной оси.

\field{Проверка.} За $4$ с скорость увеличилась на $2\cdot4=8\,\text{м/с}$. От значения $-6\,\text{м/с}$ это приводит к $2\,\text{м/с}$, что согласуется с расчётом.

\textcolor{teal!60!black}{\textbf{Итог:}} через $4$ с скорость равна $2\,\text{м/с}$, тело движется в положительном направлении оси.

\field{Задача на закрепление.} Тело имеет начальную скорость $-5\,\text{м/с}$ и ускорение $3\,\text{м/с}^2$. Найдите скорость через $3$ с и укажите направление движения в этот момент.

\newpage
\phantomsection\label{task:14-3}
\taskhead{14.09.2026 --- задание №3}

\field{Текст задания.} Конечная скорость тела равна $14\,\text{м/с}$, ускорение равно $-2\,\text{м/с}^2$, время движения --- $6$ с. Найдите начальную скорость.

\field{Ваше решение.} В исходных данных ускорение записано как положительное: $a=2\,\text{м/с}^2$. Далее использовано
\[
v_0=v-at=14-2\cdot6=2\,\text{м/с}.
\]

\field{Ваш ответ.} $v_0=2\,\text{м/с}$.

\textcolor{red}{\textbf{Ошибка:}} потерян знак минус у ускорения.

\field{Где именно возникает ошибка.} Последний корректный фрагмент --- запись конечной скорости $v=14\,\text{м/с}$ и времени $t=6$ с. Первый неверный шаг --- перенос из условия ускорения как $a=2\,\text{м/с}^2$ вместо $a=-2\,\text{м/с}^2$.

\field{Почему это неверно.} Знак ускорения является частью исходных данных и определяет, как изменяется скорость. В формуле $v=v_0+at$ отрицательное ускорение даёт отрицательное слагаемое $at$. Если заменить $-2$ на $+2$, изменится не только знак одного промежуточного действия, но и физический смысл: вместо уменьшения скорости на $12\,\text{м/с}$ получается её увеличение на $12\,\text{м/с}$. Поэтому дальнейший ответ $2\,\text{м/с}$ относится уже к другой задаче с положительным ускорением.

\textcolor{blue!60!black}{\textbf{Ключевая идея:}} при подстановке в кинематическую формулу знак ускорения нельзя отделять от его числового значения. Нужно подставлять $a=-2\,\text{м/с}^2$ целиком.

\field{Исправление.} Формулу можно сохранить:
\[
v_0=v-at.
\]
Но подставить нужно отрицательное ускорение:
\[
v_0=14-(-2)\cdot6=14+12=26\,\text{м/с}.
\]

\field{Полное правильное решение.} Используем закон скорости равноускоренного движения:
\[
v=v_0+at.
\]
Выражаем начальную скорость:
\[
v_0=v-at.
\]
Подставляем $v=14\,\text{м/с}$, $a=-2\,\text{м/с}^2$, $t=6\,\text{с}$:
\[
v_0=14-(-2)\cdot6=14+12=26\,\text{м/с}.
\]

\field{Проверка.} Подставим найденное значение в исходную формулу:
\[
26+(-2)\cdot6=26-12=14\,\text{м/с}.
\]
Получена указанная в условии конечная скорость.

\textcolor{teal!60!black}{\textbf{Итог:}} начальная скорость равна $26\,\text{м/с}$.

\field{Задача на закрепление.} Конечная скорость тела равна $11\,\text{м/с}$, ускорение равно $-3\,\text{м/с}^2$, время движения --- $4$ с. Найдите начальную скорость.

\newpage
\phantomsection\label{task:14-4}
\taskhead{14.09.2026 --- задание №4}

\field{Текст задания.} Скорость задана формулой $v_x=7-t$. Найдите начальную скорость, ускорение, момент остановки и опишите движение до остановки и после неё.

\field{Ваше решение.} Записано: $v_{0x}=7\,\text{м/с}$, $a_x=-1\,\text{м/с}^2$, из уравнения $7-t=0$ найдено $t_{\text{ост}}=7\,\text{с}$.

\field{Ваш ответ.} $v_{0x}=7\,\text{м/с}$, $a_x=-1\,\text{м/с}^2$, $t_{\text{ост}}=7\,\text{с}$. Описание движения до и после остановки отсутствует.

\textcolor{red}{\textbf{Ошибка:}} пропущено обязательное описание направления движения до остановки и после неё.

\field{Где именно возникает ошибка.} Все числовые вычисления выполнены правильно. Первый неполный шаг появляется после нахождения $t_{\text{ост}}=7\,\text{с}$: решение заканчивается, хотя условие требует ещё проанализировать знак $v_x(t)$ по обе стороны от момента остановки.

\field{Почему это неверно.} Момент остановки делит движение на два промежутка. При $0\le t<7$ имеем $7-t>0$, поэтому скорость положительна и тело движется в положительном направлении оси. При $t>7$ имеем $7-t<0$, поэтому скорость отрицательна и тело движется в противоположном направлении. Само значение времени остановки не заменяет описание этих двух режимов движения.

\textcolor{blue!60!black}{\textbf{Ключевая идея:}} чтобы описать движение по известному закону скорости, нужно определить не только момент, когда скорость равна нулю, но и знак скорости на промежутках до и после этого момента.

\field{Исправление.} Сохраняем найденные значения $v_{0x}=7\,\text{м/с}$, $a_x=-1\,\text{м/с}^2$, $t_{\text{ост}}=7\,\text{с}$ и добавляем анализ знака: до $7$ с скорость положительна, после $7$ с --- отрицательна.

\field{Полное правильное решение.} Сравним закон
\[
v_x=7-t
\]
с формулой
\[
v_x=v_{0x}+a_xt.
\]
Отсюда
\[
v_{0x}=7\,\text{м/с},\qquad a_x=-1\,\text{м/с}^2.
\]
Остановка происходит при $v_x=0$:
\[
7-t=0,\qquad t=7\,\text{с}.
\]
Для $0\le t<7$ величина $7-t$ положительна, значит тело движется в положительном направлении оси. В момент $t=7$ с тело мгновенно останавливается. Для $t>7$ величина $7-t$ отрицательна, значит после остановки тело движется в отрицательном направлении оси.

\field{Проверка.} Например, при $t=6$ с получаем $v_x=1\,\text{м/с}>0$, а при $t=8$ с --- $v_x=-1\,\text{м/с}<0$. Это подтверждает смену направления после остановки.

\textcolor{teal!60!black}{\textbf{Итог:}} $v_{0x}=7\,\text{м/с}$, $a_x=-1\,\text{м/с}^2$, остановка происходит через $7$ с; до остановки движение направлено по положительной оси, после остановки --- по отрицательной.

\field{Задача на закрепление.} Скорость тела задана законом $v_x=9-3t$. Найдите начальную скорость, ускорение, момент остановки и опишите направление движения до остановки и после неё.

\newpage
\phantomsection\label{task:14-5}
\taskhead{14.09.2026 --- задание №5}

\field{Текст задания.} На графике $v_x(t)$ известны две точки: $(1;6)$ и $(4;0)$. Найдите ускорение. Затем определите начальную скорость и запишите закон $v_x(t)$.

\field{Ваше решение.} Записаны точки $A(1;6)$ и $B(4;0)$ и формула вида
\[
a=\frac{v-v_0}{t},
\]
после чего вычисления ускорения, начальной скорости и закона движения не завершены.

\field{Ваш ответ.} Не получен.

\textcolor{red}{\textbf{Ошибка:}} решение не доведено до результата; для двух произвольных точек использована неподходящая форма вычисления ускорения.

\field{Где именно возникает ошибка.} Правильно выписаны две точки графика. Первый проблемный шаг --- переход к формуле $a=(v-v_0)/t$ без учёта того, что первая известная точка имеет время $t=1$ с, а не $t=0$.

\field{Почему это неверно.} Ускорение при линейном законе скорости равно изменению скорости, делённому на соответствующее изменение времени. Для двух точек $(t_1;v_1)$ и $(t_2;v_2)$ нужно использовать
\[
a=\frac{v_2-v_1}{t_2-t_1}.
\]
Формула $a=(v-v_0)/t$ является частным случаем, когда одно из значений скорости --- именно начальная скорость при $t=0$. В данной задаче значение $6\,\text{м/с}$ известно при $t=1$ с, поэтому считать его $v_0$ нельзя.

\textcolor{blue!60!black}{\textbf{Ключевая идея:}} сначала находите наклон прямой $v_x(t)$ по двум точкам, то есть ускорение, а затем подставляете любую из точек в $v_x=v_{0x}+a_xt$ и восстанавливаете $v_{0x}$.

\field{Исправление.} Используем обе точки:
\[
a_x=\frac{0-6}{4-1}=\frac{-6}{3}=-2\,\text{м/с}^2.
\]
Теперь подставляем точку $(1;6)$ в $v_x=v_{0x}+a_xt$:
\[
6=v_{0x}-2\cdot1,
\]
откуда
\[
v_{0x}=8\,\text{м/с}.
\]
Закон скорости:
\[
v_x=8-2t.
\]

\field{Полное правильное решение.} Для линейного графика скорости ускорение постоянно и равно наклону прямой. По точкам $(1;6)$ и $(4;0)$:
\[
a_x=\frac{v_2-v_1}{t_2-t_1}=\frac{0-6}{4-1}=-2\,\text{м/с}^2.
\]
Записываем общий вид закона скорости:
\[
v_x=v_{0x}+a_xt=v_{0x}-2t.
\]
Используем точку $(1;6)$:
\[
6=v_{0x}-2,
\]
следовательно,
\[
v_{0x}=8\,\text{м/с}.
\]
Итак,
\[
v_x(t)=8-2t.
\]

\field{Проверка.} Для $t=1$ с:
\[
v_x(1)=8-2=6\,\text{м/с}.
\]
Для $t=4$ с:
\[
v_x(4)=8-8=0.
\]
Обе исходные точки удовлетворяют найденному закону.

\textcolor{teal!60!black}{\textbf{Итог:}} $a_x=-2\,\text{м/с}^2$, $v_{0x}=8\,\text{м/с}$, $v_x(t)=8-2t$.

\field{Задача на закрепление.} На графике $v_x(t)$ известны две точки: $(2;5)$ и $(6;-3)$. Найдите ускорение, начальную скорость и запишите закон $v_x(t)$.

\newpage
\phantomsection\label{task:15-1}
\taskhead{15.09.2026 --- задание №1}

\field{Текст задания.} Скорость тела изменилась с $-4\,\text{м/с}$ до $8\,\text{м/с}$ за $6$ с. Найдите ускорение.

\field{Ваше решение.} Начальная скорость переписана как $v_0=4\,\text{м/с}$, после чего вычислено
\[
a=\frac{8-4}{6}=\frac{2}{3}\,\text{м/с}^2.
\]

\field{Ваш ответ.} $a=\frac{2}{3}\,\text{м/с}^2$.

\textcolor{red}{\textbf{Ошибка:}} потерян отрицательный знак начальной скорости.

\field{Где именно возникает ошибка.} Первый неверный шаг возникает уже при записи данных: вместо $v_0=-4\,\text{м/с}$ записано $v_0=4\,\text{м/с}$. Все последующие вычисления используют это изменённое значение.

\field{Почему это неверно.} Ускорение определяется изменением скорости за единицу времени. Если начальная скорость отрицательна, то при переходе от $-4\,\text{м/с}$ к $8\,\text{м/с}$ изменение скорости равно не четырём, а двенадцати метрам в секунду:
\[
8-(-4)=12.
\]
Знак минус в исходной скорости показывает направление движения в начальный момент и является обязательной частью числа. Его потеря меняет величину изменения скорости и, следовательно, ускорение.

\textcolor{blue!60!black}{\textbf{Ключевая идея:}} при вычислении разности скоростей отрицательное начальное значение нужно вычитать вместе со знаком: вычитание отрицательного числа превращается в сложение.

\field{Исправление.} Сохраняем формулу
\[
a=\frac{v-v_0}{t},
\]
но подставляем $v_0=-4\,\text{м/с}$:
\[
a=\frac{8-(-4)}{6}=\frac{12}{6}=2\,\text{м/с}^2.
\]

\field{Полное правильное решение.} Изменение скорости равно
\[
\Delta v=8-(-4)=12\,\text{м/с}.
\]
За время $6$ с ускорение составляет
\[
a=\frac{\Delta v}{t}=\frac{12}{6}=2\,\text{м/с}^2.
\]

\field{Проверка.} При ускорении $2\,\text{м/с}^2$ за $6$ с скорость изменяется на $2\cdot6=12\,\text{м/с}$. Начиная с $-4\,\text{м/с}$, получаем $-4+12=8\,\text{м/с}$, то есть конечное значение совпадает с условием.

\textcolor{teal!60!black}{\textbf{Итог:}} ускорение равно $2\,\text{м/с}^2$.

\field{Задача на закрепление.} Скорость тела изменилась с $-7\,\text{м/с}$ до $5\,\text{м/с}$ за $4$ с. Найдите ускорение.

\vfill
\begin{center}
\small Лёвин Артём Александрович эксклюзивно для Володи
\end{center}

\end{document}
